% Revisione Ready
% Grammatica Ready
% Citazioni Ready

\section{Minimum Hitting Set in Ambito Online}\label{sec:minimum_hitting_set_in_ambito_online}

In numerosi scenari applicativi, assumere che l'intera istanza del problema sia disponibile prima dell'esecuzione dell'algoritmo risulta poco realistico.
Questa limitazione è alla base del contesto online: nei sistemi reali le informazioni vengono osservate con il passare del tempo e, come tali, le decisioni devono essere prese senza avere la possibilità di conoscere gli aggiornamenti futuri e dovendo impiegare il minor tempo possibile.
I problemi analizzati in ambito online non hanno possibilità di conoscere gli aggiornamenti futuri in quanto l'istanza viene rivelata in maniera incrementale.
Ogni qualvolta arrivi una sequenza di iperarchi, l'algoritmo deve riuscire a verificare in modo corretto se e come serve modificare l'Hitting Set, in base alla copertura fornita dall'Hitting Set precedente e dalla durata degli iperarchi che lo stesso sta già coprendo.
Una delle caratteristiche principali che differenzia il modello online rispetto a quello dinamico è l'irrevocabilità delle decisioni.
Una volta che un nodo viene selezionato all'interno della soluzione, questo non deve più essere rimosso nelle fasi successive.
Questo vincolo apparentemente innocuo introduce una componente asimmetrica tra le operazioni disponibili e rende necessario bilanciare la qualità della soluzione corrente con l'incertezza riguardante l'evoluzione futura dell'istanza.
A riguardo, la letteratura si distingue principalmente tra gli algoritmi deterministici e quelli randomizzati, che adesso andremo ad analizzare.

\subsection*{Algoritmi deterministici}

Con il termine algoritmi deterministici si intendono tutti quegli algoritmi che affrontano il problema del Minimum Hitting Set Online seguendo regole fisse e prive di componenti randomiche.
L'assenza di casualità rende tali approcci particolarmente semplici da analizzare dal punto di vista teorico e facilmente implementabili in sistemi reali, poiché il comportamento dell'algoritmo è prevedibile e riproducibile.
Le strategie più comuni operano secondo criteri \emph{greedy}, ovvero applicando algoritmi che prendono decisioni di ottimizzazione locale con la speranza di arrivare al contempo a una minimizzazione globale della soluzione del problema.
Attraverso lo studio di un problema duale dell'Hitting Set Problem in ambito Online, ovvero il Set Cover Problem definito in ambito Online, i ricercatori Alon, Awerbuch e Azar hanno introdotto e dimostrato che nessun algoritmo deterministico può ottenere un competitive ratio migliore di $\Omega(\log m \log n)$, dove $m$ è il numero di iperarchi e $n$ il numero di vertici dell'istanza, fornendo al contempo un algoritmo $O(\log m \log n)$-competitivo che rende tale limite quasi ottimo~\cite{alon2003}.
La letteratura definisce però anche l'esistenza di limiti teorici sulle prestazioni, espressi in termini di \emph{competitive ratio}, ottenibili dagli stessi.
Tali limiti dipendono dall'istanza utilizzata, ovvero dal numero di nodi, di iperarchi o dalla frequenza massima, scelta come riferimento per l'analisi effettuata.

\subsection*{Algoritmi randomizzati}

Per superare le barriere computazionali poste dai lower bound degli algoritmi deterministici, la letteratura introduce la randomizzazione.
Tali algoritmi probabilistici si distinguono poiché non scelgono i vertici in modo univoco, ma si basano su una distribuzione di probabilità attraverso la quale definiscono con quali nodi coprire l'iperarco corrente.
L'idea alla base consiste nello sfruttare la casualità per rendere il comportamento dell'algoritmo meno prevedibile rispetto alla sequenza di input osservata.
Il vantaggio teorico della randomizzazione non può però manifestarsi in modo uniforme su tutte le varianti del problema, ma deve dipendere fortemente dalle parametrizzazioni considerate.
Per esempio, quando l'analisi viene condotta in funzione della frequenza $d$ dell'istanza, l'utilizzo della randomizzazione produce un miglioramento sostanziale: contro il limite deterministico $\Theta(d)$ discusso in precedenza, un approccio basato su un arrotondamento (\emph{rounding}) probabilistico della soluzione frazionaria, consente di ottenere un competitive ratio atteso $O(\log d)$, portando a un miglioramento esponenziale rispetto al caso deterministico~\cite{buchbinder2009}.
Questo risultato non è però generalizzabile ad ogni formulazione del problema.
Nel caso generale dell'Online Set Cover, parametrizzato rispetto al numero di vertici $n$ e di iperarchi $m$, Feige e Korman sono riusciti a dimostrare che, sotto ragionevoli ipotesi di complessità computazionale, nessun algoritmo randomizzato polinomiale può ottenere un competitive ratio migliore di $\Omega(\log m \log n)$, pareggiando di fatto il limite deterministico in questo regime~\cite{korman2004}.
Possiamo affermare dunque che la randomizzazione non permette di garantire un vantaggio assoluto sugli algoritmi deterministici, dato che la sua efficacia è strettamente legata al parametro rispetto al quale si conduce l'analisi e alla particolare struttura combinatoria o geometrica dell'istanza considerata.
Tra le tecniche più utilizzate sono presenti approcci basati su selezione probabilistica dei candidati, campionamenti incrementali e strategie che associano pesi dinamici ai vertici dell'ipergrafo, aggiornandoli nel tempo.
Quando si tratta di soluzioni basate su algoritmi randomizzati la valutazione delle prestazioni di tali algoritmi viene effettuata considerando il \emph{competitive ratio atteso}, ossia il rapporto tra il costo medio della soluzione online e quello dell'ottimo offline calcolato sull'intera sequenza.
Poiché gli algoritmi randomizzati introducono una maggiore complessità analitica e una minore prevedibilità operativa, risulta spesso necessario bilanciare la qualità attesa della soluzione e la stabilità del comportamento nei contesti applicativi reali.

\subsection*{Considerazioni finali}

Il modello online rappresenta un'estensione naturale del problema del Minimum Hitting Set (MHS), in quanto ci permette di studiare il problema all'interno di ambienti nei quali l'istanza non è interamente disponibile e le decisioni devono essere prese in modo incrementale e irrevocabile.
All'interno dell'ambito online, a causa della complessità del problema stesso, le soluzioni proposte dagli algoritmi vengono analizzate tramite il \emph{competitive ratio}, strumento che consente di confrontare il costo della soluzione proposta con il corrispettivo risultato ottimale.